\section{The bounded null-cut Ward identity}
\label{sec:null-current}

The distinction between a modular generator and a stress-tensor flux is
essential here.  The preceding construction supplies the generators
$H_g$ without identifying their action on physical fields.  We now make
that identification for bounded profiles.  Throughout this section $g$
is real and belongs to $C_b^\infty(\R^2)$, meaning that every derivative
is bounded.  All estimates involve only finitely many of these bounds.
We write $h_{\mathbf1}=\id$ for the constant profile $1$; a numerical
subscript $\nu$ in $h_\nu$ always denotes the radial Hankel operator of
Section~\ref{sec:foundations}.

The proof has two distinct outputs. First we identify matrix elements of
$H_g$ with a stress-tensor flux and extend that identity to physical
polynomial tests. Then we use its linear dependence on $g$, together
with the spectral order already proved, to construct a commuting
transverse measure. Its quadratic moment is the input to
Section~\ref{sec:dispersion}. The identification concerns bounded
profiles throughout; the unbounded parabolic generator retains its
separate role as a comparison operator.

Put
\begin{equation}\label{N:boost-current}
 j_\alpha=vT_{\alpha v}-uT_{\alpha u},\qquad
 \partial^\alpha j_\alpha=0,
 \qquad
 \mathcal C_\epsilon(\phi,Z)
   =\ip{\phi}{Z}-\epsilon\ip{Z}{\phi}.
\end{equation}
The last expression is a real bilinear form, rather than a complex
bilinear pairing.  Complex versions below are obtained by polarization.
The stress tensor is even.  Thus only the case in which the two local
tests have equal total parity contributes; $\epsilon$ is their
interchange sign.

\subsection{Bounded local tests and differentiation}

Let $A$ and $C$ be bounded Hermitian operators, localized in compact
double cones strictly inside $W$ and its opposite, respectively.  We
first smooth them by compactly supported convolution under the
Poincar\'e group.  A slightly larger double cone still contains each
support.  Such smoothing gives bounded commutators of every prescribed
finite order with the infinitesimal Poincar\'e action, and the vacuum
vectors and their products are smooth translation vectors.  On the
opposite side one may use the twist-corrected reflected net described
in Section~\ref{sec:foundations}.

The first task is to keep both tests separated while varying them with
$H_g$. This requires containment of the conjugated algebras for both
signs of the flow parameter. The bound $|H_g|\leq\norm{g}_\infty b$
alone supplies a domain estimate and does not supply that containment.

There are $\alpha,c>0$ such that the support of $A$ lies strictly inside
the truncated cut with profile $\alpha|y|^2+c$; the reflected assertion
holds for $C$.  Write $d=\norm{g}_\infty$.  For $t>0$, the positive-profile
half-sided inclusion and the constant-shift identity give
\[
 \Ad(e^{itH_g})\cM(W)=\cM(W_{tg})
       \supset\cM(W_{td})=\Ad(e^{itdb})\cM(W).
\]
Conjugating this inclusion by $e^{-itH_g}$, and using commutation
with $b$, yields
\[
 \Ad(e^{-itH_g})\cM(W)\subset\cM(W_{-td}).
\]
For positive $t$ the same upper containment follows from $tg\geq-td$.
These are algebra inclusions obtained from the cuts; they do not
follow merely from an operator-form bound.  Since $H_g$ commutes
with the parabolic generators and $b$, conjugation of
$\cM(W_{\alpha p_0+c})$ therefore lies in
$\cM(W_{\alpha p_0+c-|s|d})$.  Choosing $|s|d<c/2$ shows that
\[
 A_s=\Ad(e^{isH_g})(A),\qquad C_s=\Ad(e^{isH_g})(C)
\]
remain in fixed, strictly separated truncated regions.  Reflection
gives the assertion for $C_s$.  No identity $H_{-g}=-H_g$ is used.
Hence $A_s$ and $C_t$ graded commute for all sufficiently small $s,t$.

Set $Z=CA\Omega$.  Since $|H_g|\leq \norm{g}_\infty b$, these smooth
vectors lie in $\Dom H_g$.  Differentiating vectors, rather than
unbounded operator products, gives
\begin{equation}\label{N:product-derivative}
 Z_{C'}=\epsilon A(iH_gC\Omega),\qquad
 Z_{A'}=C(iH_gA\Omega),\qquad
 Z_{C'}+Z_{A'}=iH_gZ.
\end{equation}
For example, $C_sA\Omega=\epsilon AC_s\Omega$ in a whole interval
about zero, so the first derivative follows from strong vector
differentiation.  The product rule then follows by subtracting and
adding $C_sA\Omega$ in $C_sA_s\Omega$.  The normal form estimate of
Lemma~\ref{F:lipschitz-form} and the bounded translation commutators give
\begin{equation}\label{N:normal-domain}
 Z_{A'},Z_{C'}\in\Dom(1+a)^{1/2}.
\end{equation}
For example, Lemma~\ref{F:lipschitz-form} puts $H_gA\Omega$ in
$\Dom(1+a)^{1/2}$. The bounded commutator of $C$ with the normal
translation generator makes $C$ bounded on $\Dom(1+a)$ with its graph
norm; interpolation with its Hilbert-space bound gives the same
property on $\Dom(1+a)^{1/2}$. This proves the assertion for $Z_{A'}$,
and the roles of $A,C$ give the other one. In particular no invariance
of $\Dom a$ under $H_g$ is being used.

These are the only normal-energy domains needed for the narrow-strip
term below.  The boost calculation uses only
$\Dom B\cap\Dom H_g$. Indeed the affine unitary relation
$e^{itB}e^{isH_g}=e^{ie^t sH_g}e^{itB}$, differentiated in $t$ on
this intersection, gives
\[
 B e^{isH_g}\xi=e^{isH_g}(B+sH_g)\xi.
\]
Thus the $H_g$ unitary group preserves this intersection and is continuous
in its $B$ graph norm. Differentiating the corresponding matrix identity,
and moving $B$ onto the fixed first vector in one difference quotient,
gives
\begin{equation}\label{N:affine-form}
 i\bigl(\ip{B\xi}{H_g\eta}-\ip{H_g\xi}{B\eta}\bigr)
   =\ip{\xi}{H_g\eta},
 \qquad \xi,\eta\in\Dom B\cap\Dom H_g.
\end{equation}
Our smoothed vacuum vectors belong to this domain. No domain for
$BH_g\eta$ is required.

Decrease $c$ and relabel it so that the fixed containing truncations
for all sufficiently small $s$ have profiles $\alpha p_0+c$.
We choose a time step whose two half-slabs remain separated from the
opposite truncated regions. The geometry can be read directly from the
cut formula. Put $R_0=1/(2\alpha)$. Then
\[
 W_{\alpha p_0+c}
 =\left\{-R_0<u<0:\quad
       v>c+\frac{|y|^2}{2(R_0+u)}\right\}.
\]
For $u=-d<0$, the reflected-cut separation inequality says that
\[
 v>-c-\frac{|y|^2}{2(R_0+d)}
\]
is spacelike to the reflected opposite truncation. The margin satisfies
\begin{equation}\label{N:slab-margin}
 d+c+\frac{|y|^2}{2(R_0+d)}
 \ge c_0\sqrt{1+d^2+|y|^2}\qquad(d\ge0)
\end{equation}
for some $c_0>0$: $d+c$ controls $R_0+d$, and the
arithmetic--geometric mean controls $|y|$. Reflection gives the
corresponding inequality for $u>0$.

Choose a positive smooth width $w(u,y)$ comparable to
$\sqrt{1+u^2+|y|^2}$, with bounded first derivatives and homogeneous of
degree one outside a compact set. If $\sigma$ is a fixed smooth step,
zero below $-1$ and one above $1$, set
$\rho=\sigma((u+v)/(\varepsilon w(u,y)))$.
For sufficiently small $\varepsilon$, its transition slab is spacelike,
and \eqref{N:slab-margin} makes its negative and positive halves
spacelike to the respective opposite truncated regions. Near $u=0$ we
also require separation from both fixed compact tests $A,C$.
Their compact normal-coordinate margins give this gap on bounded
transverse sets; at large $|y|$, their quadratic spacelike separation
dominates the slab width. Shrinking $\varepsilon$ supplies the required
central gap. This last assertion concerns the compact tests, not both
entire truncated regions.

Take $\chi\in C_c^\infty(\R\times\R^2)$, real and equal to one near
zero, and write
\[
 \chi_R(u,y)=\chi(u/R,y/R),\qquad F_R=\chi_R\,\dd\rho,
 \qquad j_F\Omega=\int F\mathbin{\cdot}j(x)\Omega\,\dd^4x.
\]
The full flux vectors satisfy
\begin{equation}\label{N:full-flux}
 \sup_{R\geq R_0}\norm{j_{F_R}\Omega}<\infty,
 \qquad j_{F_R}\Omega\rightharpoonup0\quad(R\longrightarrow\infty).
\end{equation}
Indeed conservation allows replacement of $\rho$ by an entirely scaled
step agreeing with it near $\supp\dd\chi_R$.  The difference of the
fluxes is the integral of $\dd\chi_R(\rho-\rho_R)$ against $j$ and
vanishes.  The resulting compact flux has scale degree zero, so its
vector is a unitary dilation of one fixed vector.  Its spectral mass
is driven to zero under this dilation.  The stress-tensor vacuum
channel has no zero-momentum part, proving weak convergence; the norm
is unchanged.  This argument uses neither clustering nor a norm bound
on the field.

The physical Ward identities are assumed for original polynomials.
The following two mixed functionals extend them in exactly the form
needed for differentiation. In each, the boost acts on the fixed
smoothed test.

Let $\vartheta$ be smooth, equal to one on $(-\infty,-2]$ and zero on
$[-1,\infty)$, and put
\[
 \theta_{-,\delta}(u)=\vartheta(u/\delta),\qquad
 \theta_{+,\delta}(u)=\theta_{-,\delta}(-u),\qquad
 j_\pm=j_{\theta_{\pm,\delta}F_R}.
\]

\Needspace{7\baselineskip}
\begin{lemma}[Mixed boost Ward identities]\label{N:mixed-Ward}
For the bounded smoothed opposite tests above, for sufficiently large
$R$, sufficiently small $\delta>0$, and $s$ in a neighborhood of zero,
\begin{align}\label{N:undifferentiated-Ward}
 \ip{C_s\Omega}{BA\Omega}
   &=\mathcal C_\epsilon(j_-\Omega,C_sA\Omega),\notag\\
 \ip{BC\Omega}{A_s\Omega}
   &=-\mathcal C_\epsilon(j_+\Omega,CA_s\Omega).
\end{align}
\end{lemma}

\begin{proof}
Start with a Hermitian finite algebraic sum $C_0$ of factorized original
words compactly localized in the opposite wedge. Its closed affiliated
realization was constructed in Section~\ref{sec:foundations}. The boost
Ward identity applies to its original polynomial restriction; affiliation
and opposite locality allow the compact current flux around $C_0$ to
pass the bounded operator $A$.

Choose the partition for the flux around $A$ once. Its time-zero causal
shadow is compactly contained in the positive spatial halfspace, so a
small slab neighborhood of that shadow lies compactly inside $W$ and
is spacelike to every opposite-wedge polynomial. Denote the resulting
fixed compact current vacuum vector by $\eta_A$. For each $C_0$, its
own time-zero causal shadow is a compact subset of the negative
halfspace. Thin or deform the slab only near that shadow, leaving the
$A$ partition and the homogeneous shape at infinity unchanged.
Conservation and \eqref{N:full-flux}, applied to each deformation,
transfer the opposite flux to this same $\eta_A$. The signs of the
commutator give
\[
 \ip{BC_0\Omega}{A\Omega}
 =\ip{C_0\Omega}{BA\Omega}
 =\mathcal C_\epsilon(\eta_A,C_0A\Omega).
\]
In particular the current side obeys
\[
 |\mathcal C_\epsilon(\eta_A,C_0A\Omega)|
 \le2\norm{\eta_A}\norm A\norm{C_0\Omega},
\]
since $C_0A\Omega=\epsilon AC_0\Omega$. No uniform bound on the
changing flux near $C_0$ is needed. The boost has been moved to the
fixed vector $A\Omega$ before taking any Hilbert limit.

Hermitian factorized polynomial vacuum vectors are real dense in the
opposite wedge standard subspace by the Tomita-graph argument in
Section~\ref{sec:foundations}. The displayed bounded functional therefore
extends to $C_s\Omega$. For the fixed family of truncated localizations,
the compact spatial partition can be replaced by the negative half-slab:
$C_s$ is spacelike there, and the commutator vanishes outside the fixed
causal intersection of $A$. This gives the first identity in
\eqref{N:undifferentiated-Ward}. Interchanging the fixed side gives the
second. Its minus sign comes from placing the boost on the bra; the
opposite flux contributes
$\epsilon\ip{CA_s\Omega}{j_+\Omega}-\ip{j_+\Omega}{CA_s\Omega}$.
The operators $A_s,C_s$ need only lie in their fixed truncated algebras;
compact localization of these varied operators is not required.
\end{proof}

Differentiate \eqref{N:undifferentiated-Ward} in the real parameter $s$
and add. Equation~\eqref{N:affine-form} identifies the left side as
\[
 i\bigl(\ip{BC\Omega}{H_gA\Omega}
        -\ip{H_gC\Omega}{BA\Omega}\bigr)
 =\ip{C\Omega}{H_gA\Omega}.
\]
The right side is
$\mathcal C_\epsilon(j_-\Omega,Z_{C'})
-\mathcal C_\epsilon(j_+\Omega,Z_{A'})$.
Substituting \eqref{N:product-derivative} and writing the two half-slabs
as the full slab minus its central strip gives
\begin{align}\label{N:finite-variation}
 \ip{C\Omega}{H_gA\Omega}
 &=\mathcal C_\epsilon(j_{\theta_{-,\delta}F_R}\Omega,iH_gZ)
   -\mathcal C_\epsilon(j_{F_R}\Omega,Z_{A'})\notag\\
 &\quad+
 \mathcal C_\epsilon
  (j_{(1-\theta_{-,\delta}-\theta_{+,\delta})F_R}\Omega,Z_{A'}).
\end{align}
All current entries here are actual smeared vacuum vectors. This
calculation extends a bounded vector functional; it does not postulate
a stress-tensor operator commutator with an arbitrary bounded operation.

\subsection{The three cutoff estimates}

There are three different error mechanisms in
\eqref{N:finite-variation}. The narrow strip disappears weakly after
one inverse half-power of normal energy. The normal face is compared
in norm with the multiplication operator $bg(Y)$, which yields the
physical plane flux. The distant cap is controlled after dilation by
a mass-dyad sum. Keeping these mechanisms separate specifies which
vector domains are needed and which bounds must be uniform in the
normal width.

Here and below $a=(M^2+|k|^2)/(2b)$, $b>0$.  Constants in the
one-field norm can be suppressed, since there are only finitely many
polarizations.  The two useful versions of the measure are
\begin{equation}\label{N:measure-coordinates}
 \dd a\,\dd b\,\dd^2k\quad(2ab\geq |k|^2),
 \qquad
 \frac{\dd b}{2b}\,\dd^2k\,\dd m,\qquad m=M^2\geq0.
\end{equation}
All inequalities below are componentwise inequalities in the finite
rest-spin spaces, summed at the end.

We make the dependence on test functions precise.  For a compact
smooth function $q(u,y,v)$ put
\[
 p_m(q)=\sum_{|\beta|\leq m}
  \norm{(1+|u|+|y|+|v|)^m\partial^\beta q}_{L^1}.
\]
In estimates involving a step, this seminorm is applied to its $v$
derivative and to its products with the indicated compact functions
of $(u,y)$, not to the step itself.  At a fixed $R$, all constants
denoted $C_{R,N}$ are bounded by a polynomial in the corresponding
$p_{4N+20}$ seminorms, the derivatives of $\chi_R$ on their compact
support, and the derivatives of the fixed step $\vartheta$.  After
rescaling a cap by $R$, use the same seminorms of its fixed smooth
shapes; the constants are then independent of $R$ and of the sharp
normal width.  Factors $u$, $v$, or a coordinate derivative appearing
in the current are included in these seminorms.  This convention
allows uniformity for compact families of local tests and for all
the smoothing operations used below.

First consider the last term in \eqref{N:finite-variation}.  Its
normal support has width $O(\delta)$.  Fourier transformation gives
a factor $\delta$ times a rapidly decreasing function of
$(\delta a,b,k)$; the part containing $uT_{\alpha u}$ has an additional
$\delta$.  The polarization weights for indices $u,i,v$ are
$a,\sqrt{ab},b$.  Consequently, after applying $(1+a)^{-1/2}$,
the square norm is bounded by a constant times
\begin{equation}\label{N:strip-bound}
 \int_0^\infty
  \bigl\{\delta^2(1+a)+\delta^4(1+a)^3\bigr\}
     (1+\delta a)^{-6}\,\dd a\leq C_R.
\end{equation}
The omitted $(b,k)$ integrals are convergent rapid-decay integrals;
extending their domain only enlarges this bound.  On every compact
momentum set the weighted Fourier coefficients tend to zero.
Uniform boundedness and compact-momentum density therefore give
weak convergence to zero.  By \eqref{N:normal-domain},
\begin{equation}\label{N:strip-vanishing}
 \lim_{\delta\downarrow0}
 \mathcal C_\epsilon
 (j_{(1-\theta_{-,\delta}-\theta_{+,\delta})F_R}\Omega,Z_{A'})=0.
\end{equation}

Thus the strip estimate uses exactly the half-energy domain established
before the differentiation. We turn to the first term, whose sharp
normal boundary will survive as the Ward flux.

Next integrate the first term of \eqref{N:finite-variation} by parts.
Moving $iH_g$ to the first entry means applying $-iH_g$ there.  Before
that application, conservation replaces the flux by
\begin{equation}\label{N:boundary-cap}
 B_{R,\delta}-E_{R,\delta}
 =\int(-\theta_{-,\delta}')\chi_R\rho\,j_v\Omega\,\dd^4x
 -\int\theta_{-,\delta}\rho
   \bigl[(\partial_u\chi_R)j_v-(\partial_i\chi_R)j_i\bigr]
       \Omega\,\dd^4x.
\end{equation}
For this integration first introduce a smooth $v$ cutoff of length
$S$.  At fixed $R,\delta$, let $S\to\infty$ before making any other
limit.  At $b>0$, the transforms of $\rho$ and $v\rho$ have respectively
the bounds $b^{-1}$ and $b^{-2}$ times rapid large-$b$ symbols.  This
follows by differentiating the compact function $\partial_v\rho$;
the possible distributions supported at $b=0$ are killed by the
displayed stress-tensor powers of $b$.  The error containing the
derivative of the $v$ cutoff, after multiplication by $b$, is bounded
by
\[
 C_{R,\delta,N}\,b(abS+a^2)(1+Sb+a+|k|)^{-N}.
\]
Its square integral in \eqref{N:measure-coordinates} is $O(S^{-3})$
for $N$ large: set $B=Sb$, and enlarge the cone to the full positive
$(a,b)$ quadrant.  Since $\norm{H_g\xi}\leq
\norm{g}_\infty\norm{b\xi}$, this removes the cutoff in the required
matrix element.  The remaining terms converge after multiplication
by $b$, using normal Fourier decay at fixed $\delta$.

We now estimate the normal face $B_{R,\delta}$.  Its decomposition
requires only the scalar comparison channels with $\nu=0,1$:
\[
 T_{vv}=b^2c_T,\qquad
 T_{iv}/b=M s_{zi}-k_i c_T,
\]
while conservation expresses $T_{uv}$ as a linear combination of
$(ab,|k|^2)c_T$ and $k_iM s_{zi}$.  Before a radial derivative, the
coefficient of the $vT_{vv}$ term has order one.  That of the
$uT_{uv}$ term has order $\delta(M^2+|k|^2)/b$; the mixed term
$\delta|k|M/b$ is bounded by the same expression.  Differentiating
in $m=M^2$ costs $\delta/b$, and
\begin{equation}\label{N:radial-identity}
 h_\nu\bigl(M^\nu f(m)\bigr)
 =-4M^\nu\bigl((1+\nu)f'(m)+mf''(m)\bigr).
\end{equation}
For example, a mixed $\nu=1$ coefficient has the form
$\delta k_iM f_\delta(m,b,k)/b$, up to fixed constants and finite sums.
Rescaling its normal test by $u=\delta s$ gives
\[
 |\partial_m^j f_\delta|
 \le C_{R,N}(\delta/b)^j
       (1+b+|k|+\delta a)^{-N},\qquad j=0,1,2.
\]
Its radial image is
\[
 b h_1\left(\frac{\delta k_iM}{b}f_\delta\right)
 =-4\delta k_iM\left(2\partial_m f_\delta
                              +m\partial_m^2 f_\delta\right).
\]
The potentially singular factor is now
$\delta |k|M/b\le\delta a$; also $\delta m/b\le2\delta a$.
These factors are absorbed by the rapid decay, leaving an overall
factor $\delta$. The $\nu=0$ terms follow from the same radial identity.
Consequently, for every $N$,
each coefficient of $bh_\nu B_{R,\delta}$ obeys
\begin{equation}\label{N:face-symbol}
 \abs{bh_\nu B_{R,\delta}(b,k,M)}
 \leq C_{R,N}\delta
      (1+b+|k|+\delta a)^{-N}.
\end{equation}
The spectral comparison \eqref{F:spectral-equation}, applied to
these channels, and the Taylor bound \eqref{F:taylor-equation} imply
\begin{equation}\label{N:face-error}
 \norm{(H_g-bg(Y))B_{R,\delta}}
 \leq C_R\norm{g''}_\infty\sqrt\delta.
\end{equation}
In particular the estimate has no unrecorded lower cutoff in $b$.
Indeed the square integral of the right side of
\eqref{N:face-symbol} is bounded by
\begin{align*}
 C\delta^2\int\frac{\dd b}{b}\,\dd^2k
  \int_0^\infty
   \left(1+b+|k|+\frac{\delta(m+|k|^2)}{2b}\right)^{-2N}\dd m
 &\leq C_N\delta
       \int\dd b\,\dd^2k(1+b+|k|)^{1-2N}\\
 &\leq C_N'\delta.
\end{align*}

For the principal part use $w=1+b+M^2/b$, which commutes with
$g(Y)$.  For any fixed $q>1/2$, the vectors
$w^{-q}bB_{R,\delta}$ are norm bounded and converge to their plane
limits.  To see this, the coefficients after multiplication by $b$
are bounded by
\[
 C_{R,N}b(1+\delta a)(1+b+|k|+\delta a)^{-N};
\]
the $m$ integral against $w^{-2q}$ is finite, and the remaining
$b,k$ integrals converge.  Dominated convergence applies, and the
term with $u$ tends to zero.  Since a translation-smooth $Z$ belongs
to $\Dom w^q$ for every fixed $q$, this proves convergence in the
pairing.  Fourier inversion in $k$ inserts $g(y)$.  Finally $ib$
acts as $\partial_v$ on $T_{vv}\Omega$; integration by parts leaves
\begin{equation}\label{N:plane-before-gap}
 \int\chi_R(0,y)g(y)
   (\rho(0,y,v)+v\partial_v\rho(0,y,v))
   \mathcal C_\epsilon(T_{vv}(0,y,v)\Omega,Z)\,\dd v\,\dd^2y.
\end{equation}
The term with $\partial_v\rho$ vanishes by the spacelike gap.

The face comparison has now removed the internal operator $h$ from
the surviving flux. Its error tends to zero with $\delta$ at each
fixed $R$, including the region $b\downarrow0$. To complete the
argument we need a different estimate for the distant cap: its
bound must remain valid after the normal step has become sharp.

It remains to control the cap $E_{R,\delta}$ uniformly as the normal
step becomes sharp.  Dilate its variables by $R$.  Its smooth shapes
then range over a fixed bounded set of the seminorms specified above,
and its current smearing has scale degree zero.  In these coordinates
the generator becomes
\begin{equation}\label{N:rescaled-generator}
 R^{-1}b h_{g_R},\qquad g_R(y)=g(Ry).
\end{equation}
Normal integration by parts gives a factor $(1+a)^{-1}$ from a
half-space step, and $(1+a)^{-2}$ if its coefficient contains $u$.
These bounds, as well as their logarithmic $a$ derivatives up to
order two, are uniform in the smoothing width.  For example, split
the normal integral into $|u|\leq(1+a)^{-1}$ and its complement.
The first part has the claimed length, or its square with an extra
$u$; on the second part integrate by parts.  The derivative of the
step has uniformly bounded total variation, and its support has
$|u|=O(\delta/R)$.  Repeating this argument proves the stated symbol
bounds.  Tangential integrations by parts give arbitrary rapid
decay in $b,k$, with the same $b^{-2},b^{-1}$ factors from $v\rho,\rho$.

Let $F_v,F_i$ denote the resulting cap coefficients before
multiplication by $b$.  The polarization weights give, for any $N$,
\begin{equation}\label{N:cap-symbols}
 |F_v|\leq C_N(1+a)^{-1}\tau_N(b,k),\qquad
 |F_i|\leq C_N\frac{\sqrt{a/b}}{1+a}\tau_N(b,k),\qquad
 \tau_N=(1+b+|k|)^{-N}.
\end{equation}
The same inequalities hold after up to two logarithmic $M$
derivatives, with a different $C_N$.  This last assertion can also
be checked directly in each tetrad coefficient: every occurrence
of $a$ is $(M^2+|k|^2)/(2b)$, so $M\partial_Ma\leq2a$.

Use a smooth dyadic partition $\eta_{\mathrm{lo}}(M)+\sum_{j\geq0}\eta_j(M)=1$
with $\eta_{\mathrm{lo}}$ supported in $M\leq2$ and
$\eta_j$ supported in $N_j/2\leq M\leq2N_j$, $N_j=2^j\geq1$,
and uniformly bounded logarithmic derivatives.  If $F_j=\eta_jF$
and $F_{\mathrm{lo}}=\eta_{\mathrm{lo}}F$, then
\begin{equation}\label{N:cap-dyads}
 \norm{bF_{\mathrm{lo}}}\leq C,\qquad
 \norm{bF_j}\leq C,\qquad
 \norm{bhF_j}\leq C N_j^{-2}\quad(j\geq0).
\end{equation}
Here the last expression uses the high-mass pieces only.  For clarity,
the potentially worst integral in the first two bounds is
\[
 \int_{M\asymp N}\frac{\dd m}{2b}\,
       \frac{ab}{(1+a)^2}
 =\frac12\int_{M\asymp N}\frac{a}{(1+a)^2}\,\dd m
 \leq Cb\quad(N\geq1),
\]
followed by the convergent $(b,k)$ integral with $\tau_N^2$.
For $M\leq2$ the integrand is simply bounded and the $m$ interval
has bounded length.  Thus the factor $b^{-1/2}$ in
\eqref{N:cap-symbols} is harmless.  To prove the $h$ bound, decompose
into the finite pure-label channels of
\eqref{F:comparison-table}.  On $M\asymp N_j$, the identity
\eqref{N:radial-identity} and the logarithmic derivative bounds give
a factor $N_j^{-2}$, including when derivatives hit $\eta_j$.
Apply \eqref{F:spectral-equation} to each channel and sum their
norms.  No commutation of $h$ with invariant mass is asserted or used.

Write $h_{g_R}=g_R(Y)+D_R$.  The bounded-profile estimate and
\eqref{F:taylor-equation} give simultaneously
\[
 \norm{D_R\xi}\leq2\norm{g}_\infty\norm{\xi},\qquad
 \norm{D_R\xi}\leq\tfrac12R^2\norm{g''}_\infty\norm{h\xi}.
\]
The second is an operator-norm estimate on $\Dom h$, since $D_R$
strongly commutes with $h$.  Thus the high dyads of the remainder
in \eqref{N:rescaled-generator} have norms at most
\[
 \frac{C_g}{R}\min(1,R^2/N_j^2),
\]
and their sum, including the low-mass part, is
\begin{equation}\label{N:cap-remainder}
 O_g\bigl((1+\log R)/R\bigr).
\end{equation}
Indeed, the dyads with $N_j\leq R$ contribute at most
$C_g/R$ each, and there are $O(1+\log R)$ of them. The remaining
bound is a geometric series in $R^2/N_j^2$, whose sum is bounded
independently of $R$. This gives the stated rate without any
orthogonality assumption on the images of the dyads.

This estimate allows $h$ to mix masses: the summation is a triangle
inequality after applying $D_R$, not an orthogonal decomposition
of its image.  The principal part $R^{-1}bg_R(Y)$ does preserve
mass.  Pair its $j$th dyad against the correspondingly rescaled
$Z_R$.  Translation smoothness gives, for every fixed integer $r$,
\[
 \norm{1_{M\asymp N_j}Z_R}
 \leq C_r(Z)(1+N_j/R)^{-r}.
\]
Consequently its pairing is bounded by
\begin{equation}\label{N:cap-principal}
 \frac{C_gC_r(Z)}R
 \left(1+\sum_{j\geq0}(1+N_j/R)^{-r}\right)
 \leq C_{g,r,Z}\frac{1+\log R}{R},\qquad r\geq1.
\end{equation}
Equations~\eqref{N:cap-remainder}--\eqref{N:cap-principal} prove
cap vanishing in precisely the order required by
\eqref{N:finite-variation}.

\subsection{The Ward formula and polynomial tests}

The order of limits is
\begin{equation}\label{N:limit-order}
 S\longrightarrow\infty,\qquad
 \delta\downarrow0\text{ at fixed }R,\qquad
 R\longrightarrow\infty.
\end{equation}
The strip term vanishes by \eqref{N:strip-vanishing}, the cap by
\eqref{N:cap-remainder}--\eqref{N:cap-principal}, and the full-flux
term by \eqref{N:full-flux}.  The normal face gives
\eqref{N:plane-before-gap}.  In particular its dependence on $g$
is linear before the last limit.  This proves linearity of $H_g$
as a form on these local vectors for all profiles under consideration.
The estimate
\[
 |\ip{\xi}{H_g\eta}|
 \leq\norm{g}_\infty\norm{b^{1/2}\xi}\norm{b^{1/2}\eta}
\]
extends this equality to the whole form domain.  The needed graph
density follows by applying a small compact translation convolution
to Reeh--Schlieder approximants in any fixed double cone: convolution
is bounded from Hilbert norm to the graph norm of $b^{1/2}$, and its
approximate identity tends to the identity in that graph norm.

\Needspace{6\baselineskip}
\begin{proposition}[Bounded-profile Ward formula]\label{N:Ward-proposition}
The map $g\mapsto H_g=bh_g$ is real linear on $C_b^\infty(\R^2)$.
For $g$ compactly supported or Schwartz, and the bounded smooth
opposite tests above,
\begin{equation}\label{N:bounded-Ward}
 \ip{C\Omega}{H_gA\Omega}
 =\int_{\R^2}\int_\R g(y)\rho(v)
    \mathcal C_\epsilon(T_{vv}(0,y,v)\Omega,CA\Omega)
        \,\dd v\,\dd^2y.
\end{equation}
Here $\rho$ is any smooth step whose derivative lies in the $v$ gap
between the two causal intersections with the plane.  The expression
is a vacuum-vector distribution paired with the indicated smooth
vector, not an unsmeared operator integral.
\end{proposition}

To replace the original plane step by this $\rho(v)$, interpolate
inside the spacelike gap and use locality.  The half-step vacuum
wave is continuous after an inverse polynomial energy weight:
$T_{vv}$ supplies $b^2$, the step costs $b^{-1}$, and tangential
smearing supplies rapid Fourier decay.  This is also the weighted
convergence proved in the normal-face calculation.  It justifies
the plane restriction and the removal of $\chi_R$ in
\eqref{N:bounded-Ward}.

\Needspace{6\baselineskip}
\begin{proposition}[Polynomial extension]\label{N:polynomial-extension}
Equation~\eqref{N:bounded-Ward} holds for compactly smeared physical
polynomial tests, including their adjoints, on the common field
domain $\mathcal D_{\mathrm{old}}$.  The convergence to it is uniform when the smearing tests
vary over a compact smooth parameter family.
\end{proposition}

The extension has three separate passages, with uniformity for fixed word
labels and lengths and a compact smooth family of their tests. For finite sums of factorized
polynomials, polar truncation first replaces each polynomial by a bounded
operation in its original localization algebra. Group convolution then
gives the translation derivatives needed in the Ward formula. The rapid
truncation error absorbs the polynomial growth of the norms of the other
bounded factors; only after that limit do we remove the convolution.
Finally, test-space density and the original joint-distribution bounds
extend the identity to admitted compact joint kernels.

\begin{proof}
First let $A$ be a Hermitian finite algebraic sum of factorized physical
polynomial words on the invariant field domain $\mathcal D_{\mathrm{old}}$.
Sums and products of its affiliated single-field realizations,
followed by closure, give a closed affiliated operator $T$ such that
$T|_{\mathcal D_{\mathrm{old}}}=A$ and
$T^*|_{\mathcal D_{\mathrm{old}}}=A$.  The product is
closable because its adjoint contains the reversed adjoint product
on the dense domain $\mathcal D_{\mathrm{old}}$.  Write $T=U|T|$, let
$E_L=1_{[0,L]}(|T|)$, and set
\[
 T_L=TE_L,\qquad A_L=\tfrac12(T_L+T_L^*).
\]
Then $A_L$ is bounded Hermitian in the same local algebra, with
$\norm{A_L}\leq L$.  For $\psi\in\mathcal D_{\mathrm{old}}$ and every nonnegative integer $m$,
the spectral-tail bounds give
\begin{align}\label{N:polar-tails}
 \norm{(T-T_L)\psi}
 &\leq L^{-2m}\norm{|T|^{2m+1}\psi},\notag\\
 \norm{(T^*-T_L^*)\psi}
 &\leq L^{-2m}\norm{|T|^{2m}T^*\psi}.
\end{align}
Both right sides are finite polynomial-domain norms.  Indeed
$\mathcal D_{\mathrm{old}}$ is invariant under $T,T^*$, and powers of $T^*T$ and
$TT^*$ there agree with the corresponding iterated physical
polynomials.  Hence $A_L\psi-A\psi=O(L^{-N})$ for every prescribed
$N$.  This proves the same assertion for any fixed finite family
of polynomial states and moments needed in the argument.  No
essential self-adjointness is required.

One must not differentiate the sharp spectral projection $E_L$.
Instead first convolve $A_L$ under the spacetime group, with a fixed
compact smooth kernel $f$.  Group derivatives then fall on $f$.
The estimate \eqref{N:polar-tails} is uniform for the transformed
polynomial tests on the compact support of $f$, since its right
sides are continuous norms of smeared polynomial vectors.  Thus
the convolved vacuum vectors converge in every fixed translation
Sobolev norm.  Apply the same construction to $C$ separately.
Their localization regions remain spacelike.  For example,
\[
 C_LA_L\Omega-CA\Omega
 =C_L(A_L-A)\Omega+(C_L-C)A\Omega,
\]
whose first term is bounded by $L$ times a faster-than-any-power
tail and whose second is another polynomial-state tail.  The same
argument after convolution gives convergence of products in the
finite Sobolev norms used above.  To make the product convergence explicit, let $\alpha_r$ denote the
Poincar\'e action and write $A_L(r)=\alpha_r(A_L)$,
$A(r)=\alpha_r(A)$, with analogous notation for $C$. Uniformly for
$r,s$ in fixed compact group sets,
\[
 \norm{C_L(s)A_L(r)\Omega-C(s)A(r)\Omega}
 \leq L C_NL^{-N}+C_N'L^{-N}.
\]
For the second term apply the polar-tail estimate to the polynomial
state $A(r)\Omega$; its iterated $C(s)$-polynomial norms are uniformly
bounded on these compact sets. For a product of the two convolved
operators, a common translation is left translation of both group
variables. Integration by parts consequently places each prescribed
translation derivative on the convolution kernels. The resulting
finite sum of compact kernel integrals has the same tail estimate,
with new constants. This proves the required Sobolev convergence
without differentiating either polar projection.

For these factorized polynomials, the left side of
\eqref{N:bounded-Ward} converges by its $b^{1/2}$ form bound, and
the right side by the inverse-energy estimates for the half-step
wave.  Finally remove the convolution on the original smooth field
tests.  All constants are finite test seminorms, so this last limit
and compact-parameter uniformity follow from distributional
continuity.

It remains to pass to the admitted compact joint kernels. For an old
word $A$, write $\Psi_A(\varphi_A)$ for its vacuum vector smeared with a
joint kernel $\varphi_A$. First partition a compact support into finitely
many product patches inside the required localization region. On each
patch, finite sums of tensor tests are dense in the smooth test-function
topology. Choose such kernels $\varphi_{A,j}\to\varphi_A$ with their supports
in those patches; their reversed adjoint kernels converge as well. Do the
same for the kernel $\varphi_C$ of $C$. Polarization gives the factorized
identity for arbitrary complex word tests, so it applies to these
approximants. This is a test-space approximation, not a finite-rank
requirement on the given kernels.

For every fixed nonnegative integer $k$, the original joint-distribution
bounds imply
\[
 \bigl\|(1+P^0)^k\bigl(\Psi_A(\varphi_{A,j})-\Psi_A(\varphi_A)\bigr)\bigr\|
 \longrightarrow0,
\]
and analogously for $C$ and for the concatenated old word with kernel
$\varphi_{C,j}\otimes\varphi_{A,j}\to\varphi_C\otimes\varphi_A$. Indeed, a collective
time-translation derivative of a vacuum word is the same word with the
corresponding derivative of its joint test. The doubled original vacuum
distributions therefore bound each displayed graph norm by finitely many
joint test seminorms. Positive translation spectrum makes these norms
control the $b^{1/2}$ graph norm and every fixed translation-Sobolev norm
used above. The left side of \eqref{N:bounded-Ward} is continuous by its
$b^{1/2}$ form bound; the right side is continuous by the inverse-energy
half-step bound applied to the concatenated old vacuum word. Taking this
kernel limit after the polar and convolution limits proves the stated Ward
identity for the joint kernels.

On a compact smooth parameter family the same finitely many seminorms are
uniformly bounded. The tensor approximations can be chosen uniformly in
those seminorms, using a finite cover of the compact family and a smooth
partition in its parameter. The preceding graph and flux bounds then give
the asserted compact-parameter uniformity. This final passage does not
assert a closed affiliated operator for an arbitrary joint smear.
\end{proof}

\subsection{The commuting measure and its second moment}

The analytic part of the argument has established the bounded Ward
identity and its polynomial extension. We now use the resulting
linearity to organize all bounded profiles simultaneously. Positivity
will give a measure, but normalization requires a separate tightness
argument, and its quadratic moment requires comparison with the
actual parabolic cuts.

Linearity upgrades the spectral order of bounded profiles to
commutativity.  To record the argument, put $A=h_g$, $B=h_f$ and
$\varepsilon=\norm{f}_\infty/s$.  The profile comparison of
Section~\ref{sec:foundations} gives the spectral brackets
\[
 A-\varepsilon\id\ \leq_{\rm sp}\ A+B/s
       \ \leq_{\rm sp}\ A+\varepsilon\id.
\]
If $r<t$ and $2\varepsilon<t-r$, the low $A$ subspace
$1_{(-\infty,r]}(A)\cH$ is contained in the spectral subspace of
$A+B/s$ below $r+\varepsilon$, while the high $A$ subspace
$1_{[t,\infty)}(A)\cH$ is contained in its spectral subspace above
$t-\varepsilon$.  Thus
\[
 1_{(-\infty,r]}(A)(A+B/s)1_{[t,\infty)}(A)=0.
\]
The same matrix of $A$ is zero, so the matrix of $B$ is zero.
Varying $r,t$ proves that $B$ commutes with every spectral projection
of $A$.  This argument uses the spectral-support comparison
\eqref{F:bounded-spectral-order}, not merely ordinary operator order.

\Needspace{6\baselineskip}
\begin{proposition}[Internal measure and quadratic moment]
\label{N:measure-proposition}
On the $b>0$ representation there is a normalized positive
operator-valued Borel measure $\mu$
on $\R^2$, with mutually commuting values on the internal
multiplicity space, such that
\begin{equation}\label{N:measure-representation}
 H_g=b\int_{\R^2}g(Y+z)\,\mu(\dd z),
 \qquad \mu(\R^2)=\id.
\end{equation}
Its values commute with $h$ and with the external variables $b,Y$.
Its first moment is zero as a polarized integral on
$\Dom h^{1/2}$. Its trace second moment is the increasing closed-form
limit of compactly truncated moments and defines a positive
self-adjoint operator $\kappa$ satisfying
\begin{equation}\label{N:moment-comparison}
 \int z\,\mu(\dd z)=0,\qquad
 \kappa=\int |z|^2\,\mu(\dd z)\leq h,
 \qquad K_n=b(|Y|^2+\kappa).
\end{equation}
The second-moment inequality is a quadratic-form inequality and,
since the operators strongly commute, has its usual joint
spectral meaning.  No assertion that $\mu$ is projection-valued,
or that $\kappa=h$, is made.
\end{proposition}

\begin{proof}
The map $g\mapsto h_g$ is positive and contractive in the supremum
norm, and $h_{\mathbf1}=\id$. It commutes strongly with $Y,h,b$ and is
boost independent. In the $Y$ representation write $h_g=M_{F_g}$,
where $F_g$ is an essentially bounded internal-operator-valued function.
For $\tau_tg(z)=g(z+t)$, transverse covariance gives
\[
 F_{\tau_tg}(y)=F_g(y+t)\quad\text{almost everywhere},\qquad
 \|F_g(\cdot+t)-F_g\|_{L^\infty}
 \le\|\tau_tg-g\|_\infty.
\]
The right side tends to zero. Weak-operator convolution with a smooth
approximate identity therefore converges to $F_g$ in essential-supremum
norm. The convolutions are norm continuous; their pairwise
essential-supremum bounds are consequently actual supremum bounds.
Their uniform limit selects a unique norm-continuous representative
$\widetilde F_g$. Covariance now holds pointwise, so
\[
 \mathcal L(g)=\widetilde F_g(0),\qquad
 \widetilde F_g(y)=\mathcal L(\tau_yg)
\]
defines one positive contractive internal functional. Linearity,
positivity, and the commuting relations pass to these representatives.
For inequalities involving the unbounded $h$, first restrict to a
bounded spectral band of $h$, evaluate the resulting continuous
inequality at zero, and then exhaust the bands.

The Riesz representation theorem applied to scalar matrix elements
of $\mathcal L$ on compact tests, and polarized, gives a positive
operator-valued measure $\mu$. Commutativity of all $h_g$ and
commutation with $h$ pass to its Borel values by bounded monotone
approximation. This proves \eqref{N:measure-representation} first for
compactly supported smooth profiles.

It remains to rule out mass at infinity, which is not automatic
from a representation on compact tests.  Choose a radial
$\chi_0\in C_c^\infty(\R^2)$, $0\leq\chi_0\leq1$, equal to one
near zero and nonincreasing along rays.  Set
$\chi_{0,R}(z)=\chi_0(z/R)$.  The Taylor estimate at the $Y=0$
fiber gives, on $\Dom h^{1/2}$,
\begin{equation}\label{N:tightness}
 0\leq\id-\mu(\chi_{0,R})
 \leq C R^{-2}h.
\end{equation}
For $\mathcal E_R=\id-\mu(\chi_{0,R})$, positivity gives
$0\leq \mathcal E_R\leq\id$, hence
$\norm{\mathcal E_R\psi}^2\leq\ip{\psi}{\mathcal E_R\psi}$. Equation~\eqref{N:tightness}
therefore gives strong convergence to zero first on
$\Dom h^{1/2}$ and then, by density and the contraction bound,
on the whole internal space.
Monotone convergence and density imply
$\mu(\chi_{0,R})\uparrow\id$ strongly.  This proves normalization
and excludes escaping mass.  It also extends the representation
from compact tests to bounded smooth profiles.  Indeed positivity
and linearity at the $Y=0$ fiber give
\[
 |\mathcal L(g)-\mu(g\chi_{0,R})|
 \leq\norm{g}_\infty\bigl(\id-\mu(\chi_{0,R})\bigr).
\]
The right side tends strongly to zero; it controls the omitted
scalar matrix elements, and polarization completes the extension.

For the precise second-moment constant, put
$q_R(z)=|z|^2\chi_{0,R}(z)$.  These bounded smooth profiles increase
to $|z|^2$ and satisfy $q_R\leq p_0$.  Direct cut inclusion and log
order give $h_{q_R}\leq |Y|^2+h$.  At $Y=0$ this is
\[
 0\leq\mu(q_R)\leq h.
\]
Monotone convergence defines the densely defined closed second
moment form $\kappa$, dominated by $h$.  All its bounded spectral
approximants commute with $h$, so $\kappa$ is affiliated with the
same commutative algebra.  This proves $\kappa\leq h$ with coefficient
one, rather than a Hessian-dependent constant.

Finally let $e\in\R^2$ and
$a_R(z)=(e\cdot z)\chi_{0,R}(z)$.  Then $a_R(0)=0$ and
$\norm{a_R''}_\infty\leq C|e|/R$, so
\[
 |\mu(a_R)|\leq C|e|R^{-1}h.
\]
The second-moment bound makes the first moment integrable on the
$h$ form domain: for $\psi\in\Dom h^{1/2}$,
\[
 \int|z|\,\ip{\psi}{\mu(\dd z)\psi}
 \leq\norm{\psi}\,\norm{h^{1/2}\psi}.
\]
Dominated convergence and polarization now give
$\int e\cdot z\,\mu(\dd z)=0$.  This proves
\eqref{N:moment-comparison}.  Covariance of the construction under
changes of transverse axes and dilations follows from covariance
of the original cut generators; boost independence has already
removed the choice of boost coordinate.
\end{proof}

The measure and its moment have been constructed using one horizon.
For the comparison between two horizons we also need them to be
independent of whether the bounded tests were taken in the original
or the reflected net. The physical Ward formula supplies precisely
that identification.

\Needspace{6\baselineskip}
\begin{corollary}[Reflected agreement]\label{N:reflected-agreement}
Repeat the construction with past cuts of the original physical
net, or with the twist-corrected reflected net, and reflect back
to the same future orientation.  The bounded profile functionals,
the measure $\mu$, and the quadratic moment $K_n$ agree.
\end{corollary}

\begin{proof}
Wightman wedge reflection gives the same polynomial vacuum data
and the same physical stress-tensor flux on the right side of
\eqref{N:bounded-Ward}.  Proposition~\ref{N:polynomial-extension}
therefore identifies the two bounded-profile forms on opposite
polynomial vacuum vectors.  Their $b^{1/2}$ graph density and the
uniform form bound identify the operators.  The compact-profile
measure and its monotone moments are consequently identical.
This argument identifies the bounded functionals and their
quadratic moment; it does not identify the full parabolic modular
generator with $K_n$ at this stage.
\end{proof}
